\documentclass[dvipsnames,handout]{beamer} % Use handout option to remove pauses.
\usetheme{default}
\usefonttheme{structurebold}
\usefonttheme[onlymath]{serif}
\usepackage[UKenglish,cleanlook]{isodate}                                  % Set default date and date display
\usepackage[T1]{fontenc}
% Slide title background color
\definecolor{background}{HTML}{ede6d8}
% Slide title text color
\definecolor{titleText}{HTML}{B40404}
% Add slide numbers in bottom right corner
\defbeamertemplate{headline}{my header}{
    \vskip1pt
    \makebox[0pt][l]{\,\insertsection}
    \hspace*{\fill}\insertshorttitle\hspace*{\fill}
    \llap{\insertframenumber\,/\,\inserttotalframenumber\,}
}
\setbeamertemplate{headline}[my header]
% Add name at the bottom
\setbeamercolor{footlinecolor}{fg=black,bg=background}
\defbeamertemplate{footline}{my footer}{%
    \begin{beamercolorbox}[wd=\paperwidth,leftskip=0.25cm,rightskip=0.25cm]{footlinecolor}
        \insertshortauthor
        \hfill
        \insertframenumber{} / \inserttotalframenumber
    \end{beamercolorbox}
}
\setbeamertemplate{footline}[my footer]
\setbeamertemplate{navigation symbols}[default]
\usepackage{graphicx}
% Set font sizes for frame title and subtitle
\setbeamerfont{frametitle}{size=\fontsize{15}{16}}
\setbeamerfont{framesubtitle}{size=\small}
% Set left and right text margins
\setbeamersize{text margin left=5mm, text margin right=5mm}
\usepackage{booktabs,multirow}
\usepackage{subcaption}   % Sub figures.
\usepackage{tikz}
\usetikzlibrary{shapes,decorations,decorations.pathreplacing,arrows,calc,arrows.meta,fit,positioning}
\tikzset{
    auto,node distance =1 cm and 1 cm,semithick,
    state/.style ={ellipse, draw, minimum width = 0.7 cm},
    point/.style = {circle, draw, inner sep=0.04cm,fill,node contents={}},
    bidirected/.style={Latex-Latex,dashed},
    el/.style = {inner sep=2pt, align=left, sloped}
}
\usepackage{mathtools}                                                     % Various maths functions
\usepackage{amssymb}                                                       % Various maths functions
\usepackage{amsmath}                                                       % Various maths functions
\usepackage{dsfont}                                                        % Various maths functions
\usepackage{centernot}                                                     % center \not usage
\usepackage{siunitx} \sisetup{round-mode=places, round-precision=3}        % Formalise use of units and numbers among text
\usepackage[normalem]{ulem} % Strike-through package
\renewcommand{\vec}[1]{\boldsymbol{\mathit{#1}}}                           % vector notation shortcut
\newcommand{\mat}[1]{\boldsymbol{\mathit{#1}}}                             % matrix notation shortcut
\DeclarePairedDelimiter\abs{\lvert}{\rvert}                                % absolute value notation shortcut
\DeclarePairedDelimiter\norm{\lVert}{\rVert}                               % norm notation shortcut
\newcommand{\Prob}[1]{\Pr\left( #1 \right)}                         % SHortcut for probability notation
\newcommand{\Probgiven}[2]{\Pr\left( #1 \, \middle\vert \, #2 \right)} % SHortcut for probability notation, given
\newcommand{\E}[2][]{\mathbb{E}_{#1} \left[ #2 \right]}                    % Expectation (with optional subscript) shortcut
\newcommand{\Egiven}[3][]{\mathbb{E}_{#1} \left[ #2 \, \middle\vert \, #3 \right]} % Expectation given (with optional subscript) shortcut
\newcommand{\Var}[2][]{\text{Var}_{#1} \left( #2 \right)}                  % Variation (with optional subscript) shortcut
\newcommand{\Cov}[1]{\text{Cov} \left( #1 \right)}                         % Covariance (with optional subscript) shortcut
\newcommand{\indicator}[1]{\mathds{1}\left\{ #1 \right\}}                  % SHortcut for indicator function
\newcommand{\indep}{\, \raisebox{0.05em}{\rotatebox[origin=c]{90}{$\models$}} \,}% Statistical independence symbol.
\newcommand{\diff}[2][]{\frac{d#1}{d#2}}                                   % SHortcut for differential fraction as a function
\newcommand{\partialdiff}[2][]{\frac{\partial#1}{\partial#2}}              % SHortcut for partial differential fraction as a function
\renewcommand{\hat}[1]{\widehat{#1}}                                       % Default estimator notation is widehat
\renewcommand{\bar}[1]{\overline{#1}}                                      % Make over bar look nicer
\renewcommand{\tilde}[1]{\widetilde{#1}}                                   % Make over tilde look better
% Citations
\usepackage{natbib}                                        % Citation package, see https://en.wikibooks.org/wiki/LaTeX/Bibliography_Management#Natbib
\usepackage{hyperref}                                        % Allow for links across the text, with colour options
\usepackage{setspace}
\settowidth{\leftmargini}{\usebeamertemplate{itemize item}}
\addtolength{\leftmargini}{\labelsep}
\usepackage{soul,color,xcolor} % Text highlighting
\newcommand{\eqhighlight}[2]{\colorbox{#1!50}{$\displaystyle#2$}}
\makeatletter
\let\HL\hl
\renewcommand\hl{%
    \let\set@color\beamerorig@set@color
    \let\reset@color\beamerorig@reset@color
    \HL}
\makeatother
\newcommand{\mathcolorbox}[2]{\colorbox{#1}{$\displaystyle #2$}}
% Set colors
\setbeamercolor{block title}{use=structure,fg=white,bg=structure.fg!75!black}
\setbeamercolor{block body}{parent=normal text,use=block title,bg=block title.bg!10!bg}
\setbeamercovered{transparent}
\setbeamercolor{postit}{fg=black, bg=yellow}
\setbeamercolor{frametitle}{bg=background, fg=titleText}
\setbeamercolor{subtitle}{fg=titleText}
% Command to align text
\renewcommand{\raggedright}{\leftskip=0pt \rightskip=0pt plus 0cm}
% Remove the useless buttons.
\setbeamertemplate{navigation symbols}{}
\useoutertheme[footline=empty,subsection=false]{miniframes}
\useinnertheme{circles}

%-------------------------------------------------------------------------------
% Title Page
%-------------------------------------------------------------------------------
\title{\color{titleText}
    \href{https://raw.githubusercontent.com/shoganhennessy/mediation-natural-experiment/main/mediation-natural-experiment-2025.pdf}{Causal Mediation in Natural Experiments}
}
\author[Senan Hogan-Hennessy, Cornell University]{
    Senan Hogan-Hennessy \\
    Economics Department, Cornell University \\ %\vspace{0.5cm}
    \href{mailto:seh325@cornell.edu}{\textcolor{blue}{seh325@cornell.edu}}
}
\date{} % Date, can be changed to a custom date

%-------------------------------------------------------------------------------
% Opening Slides
\begin{document}
% Justify text through-out.
\raggedright
%-------------------------------------------------------------------------------
\begin{frame}
    \frametitle{Notes for long version}
    \begin{enumerate}
        \item Translate the wording for everyone (mechanism, quasi-random), and be clearer about suggestive.  Use words like necessary but not sufficient.
        \item Needs a clearer introduction, which accurately overviews and previews the approach findings
        \item Novelty needs to be loud, so put it first, Write ATE $=$ ADE $+$ AIE for Levon, and enumerates folks theorems for why CM did not take off in econ but did in medicine epi psych)
        \item Evan Riehl recommends a slide with quotes from top 5s that investigates mechanisms (note the approach is necessary but not sufficient for mechanism analysis)
        \item Mention Kwon Roth result on my data, reject null then move on....
        \item Longer presentation needs clear reasoning on the IV.
    \end{enumerate}
\end{frame}
%-------------------------------------------------------------------------------
\begin{frame}
    \frametitle{Notes for long version: empirical IV}
    \begin{enumerate}
        \item Longer explanation of the IV in Oregon for applied audience
        \item Options for the included IV, mainly to consider as illustrative (and do not want people to expect a super clean IV, but then get an illustrative one).
        \item Talk through the quasi-experimental concerns (why is $D_i$ endogenous?)
        \item Show the IV set-up (clean pre-$Z_i$ first-stage, but exclusion restriction maybe lacking).  Can it be made binary to simplify the interpretation (and linearise the estimation?).
        \item Develop at least one slide that talks through the controlling for \textit{already diagnosed} illnesses
        \item See what the CM estimates look like without controlling for them already.
    \end{enumerate}
\end{frame}
%-------------------------------------------------------------------------------
%% Title page
\begin{frame}[noframenumbering, plain]
    % Print the title page as the first slide
    \vspace{1.5cm}
    \titlepage
    \begin{center}
        \vspace{-1.5cm}
        \includegraphics[width=2cm]{presentation-files/cornell}

        %\vspace{0.5cm}
        \vskip0pt plus 1filll
        \par\noindent\rule{\textwidth}{0.4pt}
        Cornell, Labor Economics Seminar \\
        17 November 2025
    \end{center}
\end{frame}

%-------------------------------------------------------------------------------
\begin{frame}
    \frametitle{Presentation Plan}
    \textbf{Introduction} \\
    First ten minutes: introduction, road-map, and preview of findings.
    \begin{itemize}
        \item Mechanism are important (show top 5 quotes), and we do not know much about them and given (at best) suggestive evidence.
        \item CM is a framework from elsewhere which gives sufficient evidence, though has not taken off in economics
        \item I give explicit reasoning for why conventional CM methods are unlikely to work in applied econ settings, then develop a structural approach to get back to what you want.
        \item Apply these methods to Oregon, showing how these methods work in practice (suggestive, conventional CM, then my MTE approach).
    \end{itemize}
    For this part, no model and no maths notation, just vibes.

    \textbf{1. Oregon} \\
    Second ten-minutes, the model and suggestive evidence in the OHIE.
    \begin{itemize}
        \item Introduce the OHIE, and say why the mechanisms are important and unknown
        \item Show suggestive evidence
        \item Question whether suggestive evidence is enough, when $D_i \to Y_i$ correlation is zero. 
    \end{itemize}

    \textbf{2. Causal Mediation (CM)} \\
    15-minutes, the CM framework. 
    \begin{itemize}
        \item Model: Define the ADE and AIE, identifying assumption, and thus sufficient evidence on this mechanism.
        \item Explain why this does not hold (e.g., Roy model), and get selection bias (figure of simulation)
        \item Show the figure of conventional CM (i.e., assume the mechanism) estimate for OHIE.  Explain why MI would not hold for Oregon
    \end{itemize}

    \textbf{3. CM with Selection}
    15-minutes, My approach to CM.
    \begin{itemize}
        \item Write the two-stage equation, giving intuition on where non-identification comes from (second-stage).
        \item Use my 3 assumption approach to identify the mediator mechanism MTE, then ADE $+$ AIE (write equations, and short guide, but be careful)
        \item Explain estimation with my simulation figures.
    \end{itemize}

    \textbf{4. Return to Oregon}
    10-minutes, my application of CM.
    \begin{itemize}
        \item Re-introduce the setting (good lottery first-stage, but hospital visits are free choice).
        \item Walk through my MTE model, with monotonicity $+$ relevance, then walk through the IV (if illustrative)
        \item Show my estimates of CM effects, then give discussion of findings.  Include an adjusted slide of not controlling for pre-diagnosed illnesses, alluding to Oster (2019) reasoning on controlling for more ailments to get a sharper answer.
    \end{itemize}
    Finding notes: very different CM effects from conventional, noting direct (psychological) effects of socialised healthcare access.
    Agreed with suggestive evidence, giving clear uncertainty on the proportion of the mechanism.
    Note $\to$ apply the Kitagawa (2015) Kwon Roth (2024) test to these data rejects null (there are direct effects). 

    \textbf{Conclusion}
    Bullet-point overview of everything that has happened..
    
\end{frame}
%-------------------------------------------------------------------------------
\section{Introduction}
\begin{frame}
    \frametitle{Causal Mediation in Natural Experiments}
    Natural experiments gives credible estimates of causal effects
    \begin{itemize}
        \item Little information on the \textcolor{ForestGreen}{mechanisms} through which they operate
        \item Limits understanding of the decisions and underlying economic system
        \item  
    \end{itemize}

    School meals $\to$ adult income, because increased education 
    (Lundborg Rooth Alex-Petersen 2022).

    \textbf{Senan note:} unsure how to start the presentation with 5 minutes that grab attention.
    Just start section 1, introducing the data succinctly.
\end{frame}
%-------------------------------------------------------------------------------
\section{1. Oregon}
\begin{frame}
    \frametitle{Oregon Health Insurance Experiment}
\end{frame}
%-------------------------------------------------------------------------------
\section{2. Causal Mediation (CM)}
\begin{frame}
    \frametitle{Causal Mediation (CM)}
\end{frame}
%-------------------------------------------------------------------------------
\section{3. CM with Selection}
\begin{frame}
    \frametitle{CM with Selection}
\end{frame}
%-------------------------------------------------------------------------------
\section{4. Return to Oregon}
\begin{frame}
    \frametitle{Oregon Health Insurance Experiment}
\end{frame}
%-------------------------------------------------------------------------------
\section{Conclusion}
\begin{frame}
    \frametitle{Conclusion}
\end{frame}
%-------------------------------------------------------------------------------
\end{document}
